% Options for packages loaded elsewhere
\PassOptionsToPackage{unicode}{hyperref}
\PassOptionsToPackage{hyphens}{url}
%
\documentclass[
  12 pt,
  a4paper,
]{article}
\usepackage{amsmath,amssymb}
\usepackage{setspace}
\usepackage{iftex}
\ifPDFTeX
  \usepackage[T1]{fontenc}
  \usepackage[utf8]{inputenc}
  \usepackage{textcomp} % provide euro and other symbols
\else % if luatex or xetex
  \usepackage{unicode-math} % this also loads fontspec
  \defaultfontfeatures{Scale=MatchLowercase}
  \defaultfontfeatures[\rmfamily]{Ligatures=TeX,Scale=1}
\fi
\usepackage{lmodern}
\ifPDFTeX\else
  % xetex/luatex font selection
  \setmainfont[]{Times New Roman}
\fi
% Use upquote if available, for straight quotes in verbatim environments
\IfFileExists{upquote.sty}{\usepackage{upquote}}{}
\IfFileExists{microtype.sty}{% use microtype if available
  \usepackage[]{microtype}
  \UseMicrotypeSet[protrusion]{basicmath} % disable protrusion for tt fonts
}{}
\makeatletter
\@ifundefined{KOMAClassName}{% if non-KOMA class
  \IfFileExists{parskip.sty}{%
    \usepackage{parskip}
  }{% else
    \setlength{\parindent}{0pt}
    \setlength{\parskip}{6pt plus 2pt minus 1pt}}
}{% if KOMA class
  \KOMAoptions{parskip=half}}
\makeatother
\usepackage{xcolor}
\usepackage[margin=1in]{geometry}
\usepackage{color}
\usepackage{fancyvrb}
\newcommand{\VerbBar}{|}
\newcommand{\VERB}{\Verb[commandchars=\\\{\}]}
\DefineVerbatimEnvironment{Highlighting}{Verbatim}{commandchars=\\\{\}}
% Add ',fontsize=\small' for more characters per line
\usepackage{framed}
\definecolor{shadecolor}{RGB}{248,248,248}
\newenvironment{Shaded}{\begin{snugshade}}{\end{snugshade}}
\newcommand{\AlertTok}[1]{\textcolor[rgb]{0.94,0.16,0.16}{#1}}
\newcommand{\AnnotationTok}[1]{\textcolor[rgb]{0.56,0.35,0.01}{\textbf{\textit{#1}}}}
\newcommand{\AttributeTok}[1]{\textcolor[rgb]{0.13,0.29,0.53}{#1}}
\newcommand{\BaseNTok}[1]{\textcolor[rgb]{0.00,0.00,0.81}{#1}}
\newcommand{\BuiltInTok}[1]{#1}
\newcommand{\CharTok}[1]{\textcolor[rgb]{0.31,0.60,0.02}{#1}}
\newcommand{\CommentTok}[1]{\textcolor[rgb]{0.56,0.35,0.01}{\textit{#1}}}
\newcommand{\CommentVarTok}[1]{\textcolor[rgb]{0.56,0.35,0.01}{\textbf{\textit{#1}}}}
\newcommand{\ConstantTok}[1]{\textcolor[rgb]{0.56,0.35,0.01}{#1}}
\newcommand{\ControlFlowTok}[1]{\textcolor[rgb]{0.13,0.29,0.53}{\textbf{#1}}}
\newcommand{\DataTypeTok}[1]{\textcolor[rgb]{0.13,0.29,0.53}{#1}}
\newcommand{\DecValTok}[1]{\textcolor[rgb]{0.00,0.00,0.81}{#1}}
\newcommand{\DocumentationTok}[1]{\textcolor[rgb]{0.56,0.35,0.01}{\textbf{\textit{#1}}}}
\newcommand{\ErrorTok}[1]{\textcolor[rgb]{0.64,0.00,0.00}{\textbf{#1}}}
\newcommand{\ExtensionTok}[1]{#1}
\newcommand{\FloatTok}[1]{\textcolor[rgb]{0.00,0.00,0.81}{#1}}
\newcommand{\FunctionTok}[1]{\textcolor[rgb]{0.13,0.29,0.53}{\textbf{#1}}}
\newcommand{\ImportTok}[1]{#1}
\newcommand{\InformationTok}[1]{\textcolor[rgb]{0.56,0.35,0.01}{\textbf{\textit{#1}}}}
\newcommand{\KeywordTok}[1]{\textcolor[rgb]{0.13,0.29,0.53}{\textbf{#1}}}
\newcommand{\NormalTok}[1]{#1}
\newcommand{\OperatorTok}[1]{\textcolor[rgb]{0.81,0.36,0.00}{\textbf{#1}}}
\newcommand{\OtherTok}[1]{\textcolor[rgb]{0.56,0.35,0.01}{#1}}
\newcommand{\PreprocessorTok}[1]{\textcolor[rgb]{0.56,0.35,0.01}{\textit{#1}}}
\newcommand{\RegionMarkerTok}[1]{#1}
\newcommand{\SpecialCharTok}[1]{\textcolor[rgb]{0.81,0.36,0.00}{\textbf{#1}}}
\newcommand{\SpecialStringTok}[1]{\textcolor[rgb]{0.31,0.60,0.02}{#1}}
\newcommand{\StringTok}[1]{\textcolor[rgb]{0.31,0.60,0.02}{#1}}
\newcommand{\VariableTok}[1]{\textcolor[rgb]{0.00,0.00,0.00}{#1}}
\newcommand{\VerbatimStringTok}[1]{\textcolor[rgb]{0.31,0.60,0.02}{#1}}
\newcommand{\WarningTok}[1]{\textcolor[rgb]{0.56,0.35,0.01}{\textbf{\textit{#1}}}}
\usepackage{longtable,booktabs,array}
\usepackage{calc} % for calculating minipage widths
% Correct order of tables after \paragraph or \subparagraph
\usepackage{etoolbox}
\makeatletter
\patchcmd\longtable{\par}{\if@noskipsec\mbox{}\fi\par}{}{}
\makeatother
% Allow footnotes in longtable head/foot
\IfFileExists{footnotehyper.sty}{\usepackage{footnotehyper}}{\usepackage{footnote}}
\makesavenoteenv{longtable}
\usepackage{graphicx}
\makeatletter
\def\maxwidth{\ifdim\Gin@nat@width>\linewidth\linewidth\else\Gin@nat@width\fi}
\def\maxheight{\ifdim\Gin@nat@height>\textheight\textheight\else\Gin@nat@height\fi}
\makeatother
% Scale images if necessary, so that they will not overflow the page
% margins by default, and it is still possible to overwrite the defaults
% using explicit options in \includegraphics[width, height, ...]{}
\setkeys{Gin}{width=\maxwidth,height=\maxheight,keepaspectratio}
% Set default figure placement to htbp
\makeatletter
\def\fps@figure{htbp}
\makeatother
\setlength{\emergencystretch}{3em} % prevent overfull lines
\providecommand{\tightlist}{%
  \setlength{\itemsep}{0pt}\setlength{\parskip}{0pt}}
\setcounter{secnumdepth}{-\maxdimen} % remove section numbering
\ifLuaTeX
\usepackage[bidi=basic]{babel}
\else
\usepackage[bidi=default]{babel}
\fi
\babelprovide[main,import]{spanish}
\ifPDFTeX
\else
\babelfont{rm}[]{Times New Roman}
\fi
% get rid of language-specific shorthands (see #6817):
\let\LanguageShortHands\languageshorthands
\def\languageshorthands#1{}
\ifLuaTeX
  \usepackage{selnolig}  % disable illegal ligatures
\fi
\usepackage{bookmark}
\IfFileExists{xurl.sty}{\usepackage{xurl}}{} % add URL line breaks if available
\urlstyle{same}
\hypersetup{
  pdfauthor={Tomàs Ferrandis Moscardó},
  pdflang={es-ES},
  hidelinks,
  pdfcreator={LaTeX via pandoc}}

\title{ANÁLISIS BIVARIADO SOBRE LIBERTADES EN ITALIA\\
Encuesta Social Europea, 2020/2022. Ronda 10}
\usepackage{etoolbox}
\makeatletter
\providecommand{\subtitle}[1]{% add subtitle to \maketitle
  \apptocmd{\@title}{\par {\large #1 \par}}{}{}
}
\makeatother
\subtitle{Técnicas de Investigación en Ciencia Política I. UBU\\
Práctica 2.}
\author{Tomàs Ferrandis Moscardó}
\date{2024-05-10}

\begin{document}
\maketitle

{
\setcounter{tocdepth}{2}
\tableofcontents
}
\setstretch{1.5}
\newpage

\renewcommand\tablename{Tabla}

\section{1 INTRODUCCIÓN}\label{introducciuxf3n}

\subsection{1.1 DESCRIPCIÓN DE LA
PRÁCTICA}\label{descripciuxf3n-de-la-pruxe1ctica}

Este trabajo versará sobre las libertades en Italia a partir de diversos
\textbf{análisis bivariados} sobre datos de la Encuesta Social Europea
de los años 2020 y 2022 ( Ronda 10). Se escogerá en cada caso el test
más adecuado y, finalmente, se hará una valoración crítica a partir de
los datos analizado. Las variables que se analizarán en las siguientes
actividades corresponden al cuestionario de la encuesta.

La práctica está relacionada con los temas 5 y 6 de la asignatura
Técnicas de la Investigación en Ciencia Política I del Grado
Universitario de Ciencia Política y Gestión Pública de la Universidad de
Burgos, curso 2023-2024.

Un primer paso será asegurar que en nuestra instalación de R están
disponibles algunas funciones necesarias para abrir el fichero de datos
externo (tipo .sav), para algún cálculo estadístico o formateo de datos
en Rmd.

\begin{Shaded}
\begin{Highlighting}[]
\ControlFlowTok{if}\NormalTok{(}\SpecialCharTok{!}\FunctionTok{require}\NormalTok{(haven))\{}\FunctionTok{install.packages}\NormalTok{(}\StringTok{"haven"}\NormalTok{)\} }\CommentTok{\# para abrir fichero .sav}
\FunctionTok{library}\NormalTok{(haven)}
\ControlFlowTok{if}\NormalTok{(}\SpecialCharTok{!}\FunctionTok{require}\NormalTok{(knitr))\{}\FunctionTok{install.packages}\NormalTok{(}\StringTok{"knitr"}\NormalTok{)\} }\CommentTok{\# Act. 2. Tabla Rd}
\FunctionTok{library}\NormalTok{(knitr)}
\ControlFlowTok{if}\NormalTok{(}\SpecialCharTok{!}\FunctionTok{require}\NormalTok{(DescTools))\{}\FunctionTok{install.packages}\NormalTok{(}\StringTok{"DescTools"}\NormalTok{)\} }\CommentTok{\#actividad 3}
\FunctionTok{library}\NormalTok{(DescTools)}
\end{Highlighting}
\end{Shaded}

Se parte de un fichero .sav con los resultados de la encuesta que se
descargará o copiará en una subcarpeta DATOS del directorio donde se
esté este mismo fichero .Rmd.

\subsection{1.2 GESTIÓN DE FICHEROS DE
DATOS}\label{gestiuxf3n-de-ficheros-de-datos}

Tanto el fichero de los datos de la encuesta (\emph{datosparc2.sav})
como cualquier fichero de resultados se ubicarán en un subdirectorio
DATOS.

\begin{Shaded}
\begin{Highlighting}[]
\CommentTok{\# El fichero .sav debe estar en una subcarpeta DATOS}
\NormalTok{nombreFichero}\OtherTok{\textless{}{-}}\StringTok{"datosprac2.sav"}
\NormalTok{directorioTrabajo}\OtherTok{\textless{}{-}}\FunctionTok{getwd}\NormalTok{()}
\NormalTok{rutaFichero}\OtherTok{\textless{}{-}}\FunctionTok{paste}\NormalTok{(directorioTrabajo,}\StringTok{"DATOS"}\NormalTok{,nombreFichero,}\AttributeTok{sep=}\StringTok{"/"}\NormalTok{)}
\end{Highlighting}
\end{Shaded}

A partir del fichero de tipo sav (\textbf{datosprac2.sav}) de la
Encuesta Social Europea, se creará el \emph{data frame} de R. Esta será
la \textbf{base de datos} inicial con todos los datos de la encuesta.

\begin{Shaded}
\begin{Highlighting}[]
\NormalTok{df\_ronda10}\OtherTok{\textless{}{-}}\FunctionTok{read\_sav}\NormalTok{(rutaFichero)}
\end{Highlighting}
\end{Shaded}

\section{2 ACTIVIDAD 1. CONTRASTE DE MEDIAS PARA DATOS
RELACIONADOS.}\label{actividad-1.-contraste-de-medias-para-datos-relacionados.}

\subsection{2.1 DESCRIPCIÓN DE LA ACTIVIDAD
1}\label{descripciuxf3n-de-la-actividad-1}

En esta primera actividad se realizará un contraste de hipótesis para
ver si dos variables presentan diferencias estadísticamente
significativas entre si.

Las variables son:

\begin{itemize}
\tightlist
\item
  \textbf{fairelcc: Las elecciones son libres y limpias en Italia}
  (Pregunta B3 del cuestionario).
\item
  \textbf{cttresac: Los tribunales en Italia tratan a todo el mundo por
  igual} (Pregunta B18 del cuestionario).
\end{itemize}

Ambas son Variables cualitativas ordinarias y admiten un valor que va
desde 0 (``No es válida en absoluto'') a 10 (``Es completamente
válido'') además de los valores considerados como no válidos a efectos
estadísticos 77, 88 y 99 (``Rechaza responder'', ``NS'' y ``NC''
respectivamente).

\subsection{2.2 PREPARACIÓN DE LA BASE DE
DATOS}\label{preparaciuxf3n-de-la-base-de-datos}

\subsubsection{Recodificación de valores no
válidos}\label{recodificaciuxf3n-de-valores-no-vuxe1lidos}

Se comprueba si hay valores que necesitan recodificarse. Serían los
valores externos al intervalo {[}0--10{]} que según el codebook de la
encuesta podrían ser 77, 88 ó 99 ya explicados.

\begin{Shaded}
\begin{Highlighting}[]
\FunctionTok{unique}\NormalTok{(df\_ronda10}\SpecialCharTok{$}\NormalTok{fairelcc)}
\end{Highlighting}
\end{Shaded}

\begin{verbatim}
##  [1]  7  6  2  5  8  1  0  3  9  4 NA 10
\end{verbatim}

\begin{Shaded}
\begin{Highlighting}[]
\FunctionTok{unique}\NormalTok{(df\_ronda10}\SpecialCharTok{$}\NormalTok{cttresac)}
\end{Highlighting}
\end{Shaded}

\begin{verbatim}
##  [1]  1  0  2  8  3  7  6  5  9  4 NA 10
\end{verbatim}

Se comprueba que ya deben estar recodificados o almenos no aparecen. No
obstante, en caso de que no fuese así se procedería de la siguiente
manera.

En primer lugar, se duplicarán los campos para hacer las modificación
sobre estos.

\begin{Shaded}
\begin{Highlighting}[]
\NormalTok{df\_ronda10}\SpecialCharTok{$}\NormalTok{rec\_cttresac }\OtherTok{\textless{}{-}}\NormalTok{ df\_ronda10}\SpecialCharTok{$}\NormalTok{cttresac}
\NormalTok{df\_ronda10}\SpecialCharTok{$}\NormalTok{rec\_fairelecc }\OtherTok{\textless{}{-}}\NormalTok{df\_ronda10}\SpecialCharTok{$}\NormalTok{fairelcc}
\end{Highlighting}
\end{Shaded}

Después, se recodifican los valores 77, 88 y 99 como NA

\begin{Shaded}
\begin{Highlighting}[]
\NormalTok{df\_ronda10}\SpecialCharTok{$}\NormalTok{rec\_fairelecc[df\_ronda10}\SpecialCharTok{$}\NormalTok{rec\_fairelecc }\SpecialCharTok{\textgreater{}=}\DecValTok{77} \SpecialCharTok{\&} 
\NormalTok{                            df\_ronda10}\SpecialCharTok{$}\NormalTok{rec\_fairelecc }\SpecialCharTok{\textless{}=}\DecValTok{99}\NormalTok{]}\OtherTok{\textless{}{-}}\ConstantTok{NA}

\NormalTok{df\_ronda10}\SpecialCharTok{$}\NormalTok{rec\_cttresac[df\_ronda10}\SpecialCharTok{$}\NormalTok{rec\_cttresac }\SpecialCharTok{\textgreater{}=}\DecValTok{77} \SpecialCharTok{\&} 
\NormalTok{                            df\_ronda10}\SpecialCharTok{$}\NormalTok{rec\_cttresac }\SpecialCharTok{\textless{}=}\DecValTok{99}\NormalTok{]}\OtherTok{\textless{}{-}}\ConstantTok{NA}
\end{Highlighting}
\end{Shaded}

\subsection{2.3 HIPÓTESIS}\label{hipuxf3tesis}

\begin{itemize}
\tightlist
\item
  Se plantea como \textbf{hipótesis nula (H0) que ambas medias son
  iguales}.
\item
  La hipótesis alternativa (H1) sería, por lo tanto, que las medias son
  distintas.
\end{itemize}

\subsection{2.4 APLICACIÓN DE
T-STUDENT}\label{aplicaciuxf3n-de-t-student}

En estos datos en que todos los valores de ambas variables estan entre 0
y 10 podría aplicarse sobre los capos originales.

\begin{itemize}
\item
  \emph{paired = T}, indica a la función que operará con 2 muestras
  emparejadas.
\item
  \emph{alternative = `two.sized'}, indica que la hipótesis alternativa
  es que las medias no son iguales.
\item
  Se establece un \textbf{nivel de confianza} del 95\%.
\end{itemize}

\begin{Shaded}
\begin{Highlighting}[]
\FunctionTok{t.test}\NormalTok{(df\_ronda10}\SpecialCharTok{$}\NormalTok{rec\_fairelecc,df\_ronda10}\SpecialCharTok{$}\NormalTok{rec\_cttresac,}
       \AttributeTok{alternative =} \StringTok{\textquotesingle{}two.sided\textquotesingle{}}\NormalTok{,}
       \AttributeTok{conf.level =}\NormalTok{ .}\DecValTok{95}\NormalTok{,}
       \AttributeTok{paired =}\NormalTok{ T)}
\end{Highlighting}
\end{Shaded}

\begin{verbatim}
## 
##  Paired t-test
## 
## data:  df_ronda10$rec_fairelecc and df_ronda10$rec_cttresac
## t = 22.183, df = 2503, p-value < 2.2e-16
## alternative hypothesis: true mean difference is not equal to 0
## 95 percent confidence interval:
##  1.110378 1.325724
## sample estimates:
## mean difference 
##        1.218051
\end{verbatim}

\subsection{2.5 CONCLUSIÓN}\label{conclusiuxf3n}

A la vista del resultado de \emph{t\_test}, se comprueba que hay una
diferencia significativa entre las medias de aproximadamente 1,22. El
intervalo de confianza, con un nivel de confianza del 95\%, es de
1.110378 a 1.325724. La probabilidad de errar por rechazar la hipótesis
nula es prácticamente nula p=2.2x10e16; muy inferior al 0,05 que se
establece habitualmente en ciencia política como límite para
considerarlo un valor estadísticamente significativo.

Por lo que \textbf{se debe rechazar la hipótesis nula que plantea la
igualdad de las medias: las medias son distintas}.

Lo que significa que existe una percepción distinta en los encuestados
entre la libertad y limpieza del proceso electoral en Italia y el trato
igualitario de su sistema judicial.

\textbf{Dirección de la diferencia en la función de R}

El valor positivo del estadístico \emph{t-valor}, de los intervalos y de
la \emph{diferencia de medias} nos indica el orden de la diferencia.
Viendo el orden de las variables en la función R, la primera variable
``fairelecc'' tiene una media en promedio mayor que la segunda
``cttresac'', por lo que la percepción de los encuestado sobre la
limpieza de las elecciones es más postiva que la que tienen sobre la
igualdad de trato del sistema judicial.

Si se inviertese el orden de las variables en la función de R nos dará
los mismos valores absolutos de \emph{t}, de \emph{intervalos} y de
\emph{diferencia de media} pero con signo negativo (cambiaría la
dirección de la diferencia). El valor \emph{p-value} seguirá igualmente
positivo o absoluto puesto que indica la significancia estadística.

\newpage

\section{3 ACTIVIDAD 2. ASOCIACIÓN DE
VARIABLES.}\label{actividad-2.-asociaciuxf3n-de-variables.}

\subsection{3.1 DESCRIPCIÓN DE LA ACTIVIDAD
2}\label{descripciuxf3n-de-la-actividad-2}

En esta actividad se procederá a comprobar la asociación entre dos
variables cualitativas ordinales.

Las variables a analizar serán:

\begin{itemize}
\item
  \textbf{rec\_hmsfmlsh}, \emph{Me daría vergüenza que un familiar
  cercano fuese gay o lesbiana}, representa la \textbf{variable
  independiente} y, por lo tanto sus categorías corresponderán a filas
  en la tabla de contingencia (Pregunta A48 del cuestionario).
\item
  \textbf{rec\_freehms}, \emph{Los gays y lesbianas deberían tener
  libertad para vivir como quieran}, representa la \textbf{variable
  dependiente} y consecuentemente sus categorías corresponderán a
  columnas en la tabla (Pregunta A47 del cuestionario)
\end{itemize}

\subsection{3.2 PREPARAR DATOS}\label{preparar-datos}

\subsubsection{Duplicación de campos}\label{duplicaciuxf3n-de-campos}

Se duplicarán las columnas del dataframe a manipular. En estos nuevos
campos se asignan las etiquetas a sus valores correspondientes.

\begin{Shaded}
\begin{Highlighting}[]
\NormalTok{df\_ronda10}\SpecialCharTok{$}\NormalTok{rec\_hmsfmlsh}\OtherTok{\textless{}{-}}\NormalTok{df\_ronda10}\SpecialCharTok{$}\NormalTok{hmsfmlsh}
\NormalTok{df\_ronda10}\SpecialCharTok{$}\NormalTok{rec\_freehms}\OtherTok{\textless{}{-}}\NormalTok{df\_ronda10}\SpecialCharTok{$}\NormalTok{freehms}
\end{Highlighting}
\end{Shaded}

\subsubsection{Recodificación. Etiquetas de
valores}\label{recodificaciuxf3n.-etiquetas-de-valores}

Se añaden las etiquetas correspondientes a los valores en los campos
tipo \emph{factor}.

\begin{Shaded}
\begin{Highlighting}[]
\CommentTok{\# Recodifcación para que se muestre los valores de las etiquetas}
\CommentTok{\# De la variable independiente: "Me daría vergüenza que un familiar..."}
\NormalTok{ df\_ronda10}\SpecialCharTok{$}\NormalTok{rec\_hmsfmlsh }\OtherTok{\textless{}{-}} \FunctionTok{factor}\NormalTok{(df\_ronda10}\SpecialCharTok{$}\NormalTok{rec\_hmsfmlsh,}
                               \AttributeTok{levels =} \FunctionTok{c}\NormalTok{(}\DecValTok{1}\NormalTok{,}\DecValTok{2}\NormalTok{,}\DecValTok{3}\NormalTok{,}\DecValTok{4}\NormalTok{,}\DecValTok{5}\NormalTok{),}
                               \AttributeTok{labels =} \FunctionTok{c}\NormalTok{(}\StringTok{"Muy de acuerdo"}\NormalTok{, }\StringTok{"De acuerdo"}\NormalTok{, }
                                        \StringTok{"Ni de acuerdo ni en desacuerdo"}\NormalTok{,}
                                         \StringTok{"En desacuerdo"}\NormalTok{, }
                                        \StringTok{"Muy en desacuerdo"}\NormalTok{)) }

\CommentTok{\# De la variable dependiente: "Los gays y lesbianas deberían tener }
\CommentTok{\# libertad para..."}
\NormalTok{ df\_ronda10}\SpecialCharTok{$}\NormalTok{rec\_freehms }\OtherTok{\textless{}{-}} \FunctionTok{factor}\NormalTok{(df\_ronda10}\SpecialCharTok{$}\NormalTok{rec\_freehms,}
                             \AttributeTok{levels =} \FunctionTok{c}\NormalTok{(}\DecValTok{1}\NormalTok{,}\DecValTok{2}\NormalTok{,}\DecValTok{3}\NormalTok{,}\DecValTok{4}\NormalTok{,}\DecValTok{5}\NormalTok{),}
                              \AttributeTok{labels =} \FunctionTok{c}\NormalTok{(}\StringTok{"Muy de acuerdo"}\NormalTok{, }\StringTok{"De acuerdo"}\NormalTok{, }
                                        \StringTok{"Ni de acuerdo ni en desacuerdo"}\NormalTok{,}
                                        \StringTok{"En desacuerdo"}\NormalTok{, }
                                        \StringTok{"Muy en desacuerdo"}\NormalTok{)) }
\end{Highlighting}
\end{Shaded}

\subsection{3.3 HIPÓTESIS}\label{hipuxf3tesis-1}

\begin{itemize}
\tightlist
\item
  Hipótesis nula: \textbf{No hay asociación estadísticamente
  significativa entre las variables}.
\item
  Hipótesis alternativa: Hay asociación estadísticamente significativa
  entre las variables
\end{itemize}

\subsection{3.4 LA TABLA DE CONTINGENCIA (TABLA
CRUZADA)}\label{la-tabla-de-contingencia-tabla-cruzada}

Se obtiene los valores de frecuencia absolutos.

\begin{Shaded}
\begin{Highlighting}[]
\NormalTok{freqtabla }\OtherTok{\textless{}{-}} \FunctionTok{xtabs}\NormalTok{(}\SpecialCharTok{\textasciitilde{}}\NormalTok{rec\_hmsfmlsh }\SpecialCharTok{+}\NormalTok{ rec\_freehms, }\AttributeTok{data =}\NormalTok{ df\_ronda10)}
\end{Highlighting}
\end{Shaded}

Pero lo que interesa es tener en cada celda el \% de casos respecto a la
fila. Teniendo en cuenta que cada fila representa una categoría de la
variable independiente.

Esto se indica con el parámetro ``1'' en la función de R
\emph{prop.table}

\begin{Shaded}
\begin{Highlighting}[]
\NormalTok{freqtabla\_100xFila}\OtherTok{\textless{}{-}}\FunctionTok{prop.table}\NormalTok{((freqtabla),}\DecValTok{1}\NormalTok{)}\SpecialCharTok{*}\DecValTok{100}
\end{Highlighting}
\end{Shaded}

Se ve que la función no aporta los marginales (fila de totales). Para
ello R dispone de otra función, \emph{margin.table} :

\begin{Shaded}
\begin{Highlighting}[]
\NormalTok{Totales}\OtherTok{\textless{}{-}} \FunctionTok{margin.table}\NormalTok{((}\FunctionTok{prop.table}\NormalTok{(freqtabla)}\SpecialCharTok{*}\DecValTok{100}\NormalTok{), }\DecValTok{2}\NormalTok{)}
\end{Highlighting}
\end{Shaded}

Combinación de las dos matrices de R (\emph{freqtabla\_100xFila} y
\emph{Totales}) para elaborar la \textbf{Tabla de contingencia o Tabla
cruzada} mediante la función \emph{rbind}

\begin{Shaded}
\begin{Highlighting}[]
\NormalTok{tablaContingencia }\OtherTok{\textless{}{-}} \FunctionTok{round}\NormalTok{(}\FunctionTok{rbind}\NormalTok{(freqtabla\_100xFila, Totales),}\DecValTok{2}\NormalTok{)}
\FunctionTok{kable}\NormalTok{(tablaContingencia, }\AttributeTok{caption=}\StringTok{"Opiniones sobre la libertad de gays y lesbianas para vivir como quieran (Variable Dependiente en columnas) según la vergüenza de tener un familiar cercano gay o lesbiana (Variable Independiente en filas)"}\NormalTok{)}
\end{Highlighting}
\end{Shaded}

\begin{longtable}[]{@{}
  >{\raggedright\arraybackslash}p{(\columnwidth - 10\tabcolsep) * \real{0.2583}}
  >{\raggedleft\arraybackslash}p{(\columnwidth - 10\tabcolsep) * \real{0.1250}}
  >{\raggedleft\arraybackslash}p{(\columnwidth - 10\tabcolsep) * \real{0.0917}}
  >{\raggedleft\arraybackslash}p{(\columnwidth - 10\tabcolsep) * \real{0.2583}}
  >{\raggedleft\arraybackslash}p{(\columnwidth - 10\tabcolsep) * \real{0.1167}}
  >{\raggedleft\arraybackslash}p{(\columnwidth - 10\tabcolsep) * \real{0.1500}}@{}}
\caption{Opiniones sobre la libertad de gays y lesbianas para vivir como
quieran (Variable Dependiente en columnas) según la vergüenza de tener
un familiar cercano gay o lesbiana (Variable Independiente en
filas)}\tabularnewline
\toprule\noalign{}
\begin{minipage}[b]{\linewidth}\raggedright
\end{minipage} & \begin{minipage}[b]{\linewidth}\raggedleft
Muy de acuerdo
\end{minipage} & \begin{minipage}[b]{\linewidth}\raggedleft
De acuerdo
\end{minipage} & \begin{minipage}[b]{\linewidth}\raggedleft
Ni de acuerdo ni en desacuerdo
\end{minipage} & \begin{minipage}[b]{\linewidth}\raggedleft
En desacuerdo
\end{minipage} & \begin{minipage}[b]{\linewidth}\raggedleft
Muy en desacuerdo
\end{minipage} \\
\midrule\noalign{}
\endfirsthead
\toprule\noalign{}
\begin{minipage}[b]{\linewidth}\raggedright
\end{minipage} & \begin{minipage}[b]{\linewidth}\raggedleft
Muy de acuerdo
\end{minipage} & \begin{minipage}[b]{\linewidth}\raggedleft
De acuerdo
\end{minipage} & \begin{minipage}[b]{\linewidth}\raggedleft
Ni de acuerdo ni en desacuerdo
\end{minipage} & \begin{minipage}[b]{\linewidth}\raggedleft
En desacuerdo
\end{minipage} & \begin{minipage}[b]{\linewidth}\raggedleft
Muy en desacuerdo
\end{minipage} \\
\midrule\noalign{}
\endhead
\bottomrule\noalign{}
\endlastfoot
Muy de acuerdo & 22.97 & 10.81 & 8.11 & 24.32 & 33.78 \\
De acuerdo & 4.07 & 36.65 & 25.79 & 26.24 & 7.24 \\
Ni de acuerdo ni en desacuerdo & 6.92 & 40.71 & 46.05 & 4.94 & 1.38 \\
En desacuerdo & 17.39 & 67.34 & 12.26 & 2.79 & 0.22 \\
Muy en desacuerdo & 66.87 & 28.80 & 1.69 & 0.60 & 2.05 \\
Totales & 30.54 & 45.02 & 16.61 & 5.18 & 2.65 \\
\end{longtable}

\subsection{3.5 CONCLUSIONES SOBRE LA
TABLA}\label{conclusiones-sobre-la-tabla}

En una primera observación sobre los \textbf{valores marginales} o
totales se debe destacar que la mayoría de encuestados está ``De acuerdo
con las libertades de los gays y lesbianas'' (45.02\%) siendo la segunda
opción ``Muy de acuerdo'' (30,54\%), además la respuesta con menor
porcentaje es la de ``Muy en desacuerdo'' (2,65\%).

Al observar los valores de las diferentes \textbf{categorías de la
variable independiente} o filas se tiene un primer bloque con las dos
categorías de encuestados que reconocen que sentirían algún tipo de
vergüenza por tener un familiar homosexual. Aquí se ve una diferencia
notable entre ambas categorías:

\begin{itemize}
\tightlist
\item
  Por un parte los que responden estar ``Muy de acuerdo'' con este
  sentimiento de vergüenza, la respuesta mayoritaria es contraria a la
  libertad, decantándose un 33,78\% por ``Muy en desacuerdo'' y un
  24,32\% por ``En desacuerdo''. No obstante aparece un dato muy
  interesant: un importante 22,97\% que se manifiesta ``Muy de acuerdo''
  con la libertad.
\item
  Por otra parte, los que se muestran ``De acuerdo'' con la afirmación
  de vergüenza, ºsu posición mayoritaria respecto a la libertad es ``De
  acuerdo'' (36,65\%) con la libertad, un dato también interesante.
  Seguidamente optan por ``En desacuerdo'' y ``Ni de acuerdo ni en
  desacuerdo'' con cifras similares (26,24\% y 25,79\% respectivamente).
\end{itemize}

En un nivel intermedio lo configuraría la categoría más ``neutral'': los
que no están ``Ni de acuerdo ni en desacuerdo'' con el tema de la
vergüenza y que también se decantan mayoritariamente por la misma
posición neutra (central en la tabla) sobre la libertad mayoritariamente
(46,05\%). Aunque es importante señalar que su segunda opción es
posicionarse ``De acuerdo'' con las libertades de gays y lesbianas
(40,71\%), se ve una relación ``Ni de acuerdo ni en desacuerdo- Ni
desacuerdo ni desacuerdo''

En un tercer bloque estarían las dos categorías que manifiestan rechazo
a la idea de sentir vergüenza por tener un familiar homosexual. Aquí
también se ven unas relaciones interesantes.

\begin{itemize}
\item
  Los encuestados que están ``En desacuerdo'', son los que más se
  posicionan mayoritariamente también como ``De acuerdo'' con las
  libertades (67,34\%).
\item
  Con cifras similares, en el extremo, los que están ``Muy en
  desacuerdo'' con la afirmación sobre avergonzarse se declaran ``Muy de
  acuerdo con la libertad'' (66,84\%).
\end{itemize}

En definitiva, entre los que rechazan la idea de sentir vergüenza por un
familiar homosexual se ve una coincidencia ``De acuerdo-desacuerdo'' y
``Muy desacuerdo-muy de acuerdo'' en lo que respeta a las libertades de
este colectivo.

La dependencia entre las variables se deduce fácilmente al ver cómo los
valores en cada categoría de la variable dependiente varían respecto a
su marginal. De hecho no se observan coincidencias.

\subsection{3.6 CHI CUADRADO}\label{chi-cuadrado}

Este coeficiente se muestra junto a las tablas cruzadas y nos indica si
existe asociación entre la variable independiente (filas) y la
dependiente (columnas). Solo arroja un valor igual o mayor que cero. Si
vale 0 significa que no hay asociación. Si es \textgreater0, existe
asociación. Cuanto mayor, más evidencia.

Al dar valores positivos, no indica la dirección de la asociación.

\subsubsection{Ejecución de Chi cuadrado en
R.}\label{ejecuciuxf3n-de-chi-cuadrado-en-r.}

En R existe la función \emph{chisq} para el cálculo de \emph{Chi
cuadrado} con dos parametrizaciones posibles.

\begin{itemize}
\tightlist
\item
  Pasando como argumentos directamente las variables a la función
  \emph{chisq.test}
\end{itemize}

\begin{Shaded}
\begin{Highlighting}[]
\NormalTok{chi }\OtherTok{\textless{}{-}} \FunctionTok{chisq.test}\NormalTok{( df\_ronda10}\SpecialCharTok{$}\NormalTok{rec\_hmsfmlsh, df\_ronda10}\SpecialCharTok{$}\NormalTok{rec\_freehms)}
\end{Highlighting}
\end{Shaded}

\begin{verbatim}
## Warning in chisq.test(df_ronda10$rec_hmsfmlsh, df_ronda10$rec_freehms):
## L'aproximació Chi-quadrat pot ser incorrecta
\end{verbatim}

\begin{Shaded}
\begin{Highlighting}[]
\NormalTok{chi}
\end{Highlighting}
\end{Shaded}

\begin{verbatim}
## 
##  Pearson's Chi-squared test
## 
## data:  df_ronda10$rec_hmsfmlsh and df_ronda10$rec_freehms
## X-squared = 1723.8, df = 16, p-value < 2.2e-16
\end{verbatim}

\begin{itemize}
\tightlist
\item
  Pasando como argumento la tabla en \emph{chisq.test}
\end{itemize}

\begin{Shaded}
\begin{Highlighting}[]
\NormalTok{chi }\OtherTok{\textless{}{-}} \FunctionTok{chisq.test}\NormalTok{(freqtabla)}
\end{Highlighting}
\end{Shaded}

\begin{verbatim}
## Warning in chisq.test(freqtabla): L'aproximació Chi-quadrat pot ser incorrecta
\end{verbatim}

\begin{Shaded}
\begin{Highlighting}[]
\NormalTok{chi}
\end{Highlighting}
\end{Shaded}

\begin{verbatim}
## 
##  Pearson's Chi-squared test
## 
## data:  freqtabla
## X-squared = 1723.8, df = 16, p-value < 2.2e-16
\end{verbatim}

\subsection{3.7 CONCLUSIÓN}\label{conclusiuxf3n-1}

El coeficiente indica la magnitud de la asociación. El valor de 1723,8
es muy superior a 0, lo que indica una asociación fuerte. Hay una
evidencia para el rechazo de la H0 (independencia de las variables) y
además, el \emph{p-valor} es muy inferior (2.2x10¹⁶) a \emph{alpha}
(0,001) por lo que hay un nivel de confianza superior al 99,9\%.

\textbf{Se debe rechazar la hipótesis nula: las variables están
asociadas.}

No obstante, se ve que la ejecución de la función R nos avisa:

\begin{verbatim}
Warning messages:

1: In chisq.test(datos$rec_hmsfmlsh, datos$rec_freehms) :
Chi-squared approximation may be incorrect
2: In chisq.test(freqtab) : Chi-squared approximation may be incorrect
\end{verbatim}

\newpage

\section{4 ACTIVIDAD 3. RESPECTO A LA LIBERTAD DE LOS MEDIOS DE
COMUNICACIÓN}\label{actividad-3.-respecto-a-la-libertad-de-los-medios-de-comunicaciuxf3n}

\subsection{4.1 DESCRIPCIÓN DE LA ACTIVIDAD
3}\label{descripciuxf3n-de-la-actividad-3}

En esta actividad se pretende analizar si existe una dependencia entre
la autoubicación ideológica de los encuestados y su percepción sobre el
nivel de crítica que ejercen los medios de comunicación respecto al
gobierno italiano.

Para ello se analizarán las variables:

\begin{itemize}
\item
  \textbf{lrscale}, \emph{¿Dónde se colocaría usted en la escala? Escala
  de 0 a 10 (0 representa la ubicación más a la izquierda y 10 la
  ubicación más a la derecha}). Esta es la \emph{variable independiente}
  (Pregunta A39 del cuestionario)
\item
  \textbf{medcrgvc}, \emph{Los medios de comunicación en Italia pueden
  criticar al gobierno. Escala de 0 a 10 ( 0 indica No es válida en
  absolutao y 10, Es complettamente válida)}. Esta es la \emph{variable
  dependiente} (Pregunta B15 del cuestionario)
\end{itemize}

\subsection{4.2 PREPARACIÓN DE LA BASE DE
DATOS}\label{preparaciuxf3n-de-la-base-de-datos-1}

Se comprueba que no hay valores que recodificar como NA

\begin{Shaded}
\begin{Highlighting}[]
\FunctionTok{unique}\NormalTok{(df\_ronda10}\SpecialCharTok{$}\NormalTok{lrscale)}
\end{Highlighting}
\end{Shaded}

\begin{verbatim}
##  [1] NA  5  8  6  4  3 10  1  2  7  9  0
\end{verbatim}

\begin{Shaded}
\begin{Highlighting}[]
\FunctionTok{unique}\NormalTok{(df\_ronda10}\SpecialCharTok{$}\NormalTok{medcrgvc)}
\end{Highlighting}
\end{Shaded}

\begin{verbatim}
##  [1]  1  3  6  5  8 10  2  0  9  7  4 NA
\end{verbatim}

Si fuese necesario recodificar por cualquier motivo (valores no válidos,
añadir etiquetas\ldots) que no es este el caso, convendría hacerlo sobre
una duplicación de estos campos de forma análoga a lo realizado a
apartados 2.2 de las dos actividades anteriores.

\subsection{4.3 HIPÓTESIS}\label{hipuxf3tesis-2}

\begin{itemize}
\tightlist
\item
  H0, hipótesis nula: \textbf{No hay asociación estadísticamente
  significativa entre las variables}.
\item
  H1, hipótesis alternativa: Hay asociación estadísticamente
  significativa entre las variables.
\end{itemize}

\subsection{4.4 COEFICIENTE DE TAU-B DE
KENDALL}\label{coeficiente-de-tau-b-de-kendall}

Se trata del test más adecuado cuando se trata de dos variables
cualitativas ordinales que tienen el mismo número de categorías com es
este caso.

\begin{Shaded}
\begin{Highlighting}[]
\NormalTok{taub }\OtherTok{\textless{}{-}} \FunctionTok{KendallTauB}\NormalTok{(df\_ronda10}\SpecialCharTok{$}\NormalTok{lrscale, df\_ronda10}\SpecialCharTok{$}\NormalTok{medcrgvc, }
                                          \AttributeTok{conf.level =} \FloatTok{0.95}\NormalTok{)}
\NormalTok{taub}
\end{Highlighting}
\end{Shaded}

\begin{verbatim}
##        tau_b       lwr.ci       upr.ci 
## -0.009130106 -0.043788752  0.025528540
\end{verbatim}

El coeficiente \emph{tau\_b} (-0.009130106) es negativo y cercano a 0
Los valores \emph{lwr.ci} y \emph{upr.ci} (-0.043788752 0.025528540
respectivamente) son los límites inferior y superior del intervalo de
confianza del 95\%.

\subsection{4.5 CONCLUSIÓN}\label{conclusiuxf3n-2}

\textbf{Sobre le coeficiente:} Aunque el coeficiente, al ser negativo
indica que cuando aumenta una variable la otra disminuye ( dirección de
la asociación), al tratarse de un valor tan cercano a 0 (-0.009130106)
este sugiere que no hay asociación estadísticamente relevante entre las
variables.

\textbf{Sobre el intervalo de confianza:} Los valores de \emph{lwr.ci} y
\emph{upr.ci} indican que, con un 95\% de confianza, el verdadero valor
del coeficiente de Kendall se encuentra dentro de él. Este intervalo
incluye al 0 (-0.043788752 \textless0\textless{} 0.025528540) hecho que
refuerza la interpretación de que no hay una asociación estadísticamente
significativa entre las variables.

Por todo esto:

\textbf{NO se puede rechazar la hipótesis nula. No hay asociación
estadística entre las variables}

\subsection{4.6 COEFICIENTE DE SOMMERS.}\label{coeficiente-de-sommers.}

Se trata de otro test de asociación para variables cualitativas
ordinales que también nos informa de la magnitud y la dirección de la
asociación mediante un coeficiente que oscilará también entre -1 y 1.

\begin{Shaded}
\begin{Highlighting}[]
\NormalTok{delta }\OtherTok{\textless{}{-}} \FunctionTok{SomersDelta}\NormalTok{(df\_ronda10}\SpecialCharTok{$}\NormalTok{lrscale, df\_ronda10}\SpecialCharTok{$}\NormalTok{medcrgvc, }
                                           \AttributeTok{conf.level =} \FloatTok{0.95}\NormalTok{)}
\NormalTok{delta}
\end{Highlighting}
\end{Shaded}

\begin{verbatim}
##       somers       lwr.ci       upr.ci 
## -0.009052743 -0.043415457  0.025309971
\end{verbatim}

El coeficiente de \emph{Somers} (-0,009052743) cercano a 0 y negativo.

Los valores \emph{lwr.ci} y \emph{upr.ci} (-0.043415457 y 0.025309971)
son los límites inferior y superior del intervalo de confianza del 95\%.

\subsection{4.7 CONCLUSIÓN}\label{conclusiuxf3n-3}

\textbf{Sobre le coeficiente:} Aunque el coeficiente, al ser negativo
indica que cuando aumenta una variable la otra disminuye ( dirección de
la asociación), al tratarse de un valor tan cercano a 0 (-0.009052743)
este sugiere que no hay asociación estadísticamente relevante entre las
variables.

\textbf{Sobre el intervalo de confianza:} Los valores de \emph{lwr.ci} y
\emph{upr.ci} indican que, con un 95\% de confianza, el verdadero valor
del coeficiente de Sommers se encuentra dentro de él. Este intervalo
incluye al 0 (-0.043415457 \textless0 \textless0.025309971) hecho que
refuerza la interpretación de que no hay una asociación estadísticamente
significativa entre las variables.

Por todo esto:

\textbf{NO se puede rechazar la hipótesis nula. No hay asociación
estadística entre las variables}

\section{5 REFLEXIÓN SOBRE LOS
ANÁLISIS}\label{reflexiuxf3n-sobre-los-anuxe1lisis}

Las libertades vistas estan operacionalizadas como variables ordinales
pero pueden ser tratadas como cuantitativas para obterne más información
como se estudió en la anterior práctica. Así se puede ver que los
italianos encuestados ``aprueban'' su sistema electoral con un 5,8 de
media y una mediana de 6 mientras que están cerca de suspender a sus
tribunales por el trato desigual (media de 4,52 y apenas una mediana de
5). Teniendo en cuenta que se compara la confianza en el proceso un de
máxima participación y reglamentación con la práctica diaria de órganos
de funcionarios y que Italia es una democracia consolidada, la poca
confianza en la limpieza de las elecciones es un hecho preocupante.

Los encuestados mayoritariamente no se identifican con el sentimiento de
vergüenza por tener un familiar homosexual (media de 3,8 y mediana de 4
en un intervalo de 1 a 5). Un sentimiento relacionado con su
posicionamiento a favor de la igualdad de libertades de este colectivo
(una media de 2,05 y una mediana de 2 en un intervalo de 1 a 5) esto es
coherente con la dependencia entre varibles vista. Sería interesante un
seguimiento de las rondas siguientes para ver la evolución.

Con independencia de la ideología, los encuestados aprueban la
independencia respecto al gobierno de sus medios de comunicación siendo
los valores de 6 y 8 los más destacados ( entre 0 y 10). En una
valoración bastante positiva tratándose de un modelo de prensa latino.

\section{6 EJECUCIÓN DEL Rmd Y REPOSITORIO DE
FICHEROS}\label{ejecuciuxf3n-del-rmd-y-repositorio-de-ficheros}

Para la ejecución del código R, está disponible el fichero ejecutable
Rmd y el fichero de datos \emph{sav} en el repositorio GitHub. En el
mismo repositorio se encuentra este PDF y un fichero HTML renderizados a
partir del Rmd.

El fichero \textbf{datosprac2.sav} debe ubicarse en una subcarpeta del
directorio de trabajo con nombre \emph{DATA}.

\begin{longtable}[]{@{}
  >{\centering\arraybackslash}p{(\columnwidth - 4\tabcolsep) * \real{0.2500}}
  >{\raggedright\arraybackslash}p{(\columnwidth - 4\tabcolsep) * \real{0.3500}}
  >{\raggedright\arraybackslash}p{(\columnwidth - 4\tabcolsep) * \real{0.4000}}@{}}
\toprule\noalign{}
\begin{minipage}[b]{\linewidth}\centering
\includegraphics[width=0.1\textwidth,height=\textheight]{../../recursos/iconohyperlink.jpg}
\end{minipage} & \begin{minipage}[b]{\linewidth}\raggedright
\end{minipage} & \begin{minipage}[b]{\linewidth}\raggedright
Enlace a GitHub
\end{minipage} \\
\midrule\noalign{}
\endhead
\bottomrule\noalign{}
\endlastfoot
\href{https://tofermos.github.io/cienciapoliticaygestionpublica/elecciones/italia/DATOS/datosprac2.sav}{\includegraphics[width=0.1\textwidth,height=\textheight]{../../recursos/iconosav.png}}
& Fichero de datos tipo sav &
\href{https://tofermos.github.io/cienciapoliticaygestionpublica/elecciones/italia/DATOS/datosprac2.sav}{datosprac2.sav} \\
\href{https://tofermos.github.io/cienciapoliticaygestionpublica/elecciones/italia/libertadesItalia2022.Rmd}{\includegraphics[width=0.1\textwidth,height=\textheight]{../../recursos/rmarkdown.png}}
& Fichero Rmd &
\href{https://tofermos.github.io/cienciapoliticaygestionpublica/elecciones/italia/libertadesItalia2022.Rmd}{libertadesItalia2022.Rmd} \\
\href{https://tofermos.github.io/cienciapoliticaygestionpublica/elecciones/italia/libertadesItalia2022.html}{\includegraphics[width=0.1\textwidth,height=\textheight]{../../recursos/iconohtml.png}}
& Fichero HTML &
\href{https://tofermos.github.io/cienciapoliticaygestionpublica/elecciones/italia/libertadesItalia2022.html}{libertadesItalia2022.html} \\
\href{https://tofermos.github.io/cienciapoliticaygestionpublica/elecciones/italia/libertadesItalia2022.pdf}{\includegraphics[width=0.1\textwidth,height=\textheight]{../../recursos/iconopdf.png}}
& Fichero PDF &
\href{https://tofermos.github.io/cienciapoliticaygestionpublica/elecciones/italia/libertadesItalia2022.pdf}{libertadesItalia.pdf} \\
\href{https://tofermos.github.io/cienciapoliticaygestionpublica/}{\includegraphics[width=0.2\textwidth,height=\textheight]{../../recursos/iconogithub.png}}
& &
\href{https://tofermos.github.io/cienciapoliticaygestionpublica/}{Repositorio
tofermos} \\
\end{longtable}

\end{document}
